% Q14_gold_answers.tex — GRADER CALIBRATION FIXTURE (correct, paraphrased).
% A drop-in replacement for paper/Q14_answers.tex (spec §4.7, §6.12): every answer is
% correct, but the route, the wording and the examples are deliberately far from the
% reference, so the grader must recognise correctness by substance (the rubric's
% probes), not by matching text. B1 reasons from the interval 0..10 first and uses a
% witness of its own for C4 (the reference's is another); B2 words every field its own
% way and orders the facts differently.
%
% EXPECTED[B1][correctness][4]: correct; Probe A on each candidate agrees with every yes and no of (a), and Probe A' passes on every witness of (b): x = 0 for C2's (i) and (10, 10) for C4's (ii) (captured 2026-10-11: 1 passed, 0 failed, each). No finding.
% EXPECTED[B2][card-fidelity][4]: every field agrees with the pinned range (lines 2588-2631); Budget and Variance correctly say what the range does not give. No finding on any field.
% Expected per-problem grades: B1 = 4, B2 = 4 (3 is tolerated only with a reading note and no flaw).
% A calibration run FAILS if any problem scores <= 2 or if any finding claims a false fact,
% a wrong witness, a wrong instrument outcome or a contradiction with the range.
\documentclass[11pt]{article}
\input{preamble}

\title{Q14 --- Practice II-B \\ \large Written answers}
\author{learner}
\date{}

\begin{document}
\maketitle
\sloppy

% ---------------------------------------------------------------------------
\section*{Problem B1 (concept)}\label{ans:B1}
\begin{answer}
The trait admits exactly the arguments $0, 1, \dots, 10$, and on those it
promises a result strictly above the argument; both conditions are about that
interval only. \\
$C_1$: (i)? yes \quad (ii)? yes \quad witnesses: none; the interval sits
inside $x \ge 0$, and the promise is unchanged. \\
$C_2$: (i)? no \quad (ii)? yes \quad witnesses: (i) at $x = 0$, the one
argument of the interval that $1 \le x \le 10$ leaves out. \\
$C_3$: (i)? yes \quad (ii)? yes \quad witnesses: none; the precondition is
the trait's, and a result of exactly $x + 1$ is above $x$. \\
$C_4$: (i)? yes \quad (ii)? no \quad witnesses: (ii) at $(10, 10)$: $10$ is
admitted, $10 \ge 10$ holds, $10 > 10$ does not. \\
$C_5$: (i)? yes \quad (ii)? yes \quad witnesses: none; every admitted $x$
is at most $10$, and on the interval the hypothesis $x \ge 0$ is true, so the
promise is $y > x$ there. Outside the interval it promises nothing, which no
client of the trait can notice. \\
\textbf{(c)} (i) is what makes the call legal: a generic client proves only
the trait's precondition before it calls. (ii) is what makes the client's next
step sound: after the call it uses the trait's postcondition, so the result
must satisfy it.
\end{answer}

% ---------------------------------------------------------------------------
\section*{Problem B2 (paper card)}\label{ans:B2}
\begin{answer}
\textbf{Source:} cheatsheet entry 06.9, lines 2588--2631 (the pinned range).

\textbf{Task shape:} Proof repair in Isabelle: a model is handed one of 12
frozen broken proofs (half real past failures, half injected corruptions) from
four developments and must make it build, under one of five repair workflows.

\textbf{Oracle:} Two checks, both required: the theory builds, and an
independent check finds the protected text byte-identical, the target
declaration present and none of \texttt{sorry}, \texttt{oops},
\texttt{axiomatization}, \texttt{oracle} used.

\textbf{What is trusted:} The Isabelle build, plus a contract checker that is
only syntactic and itself unverified.

\textbf{Cheat surface:} Editing what was supposed to stay fixed. A green build
can hide it: 4.2\% of the Isabelle-accepted candidates (6 of 144) touched
protected text, every one from an iterative workflow over the full theory.
Limiting edits to the proof body removes this but also forbids helper lemmas.

\textbf{Metric:} Per workflow, valid repairs out of 36 and false successes:
C4 reached 32 and C1 31, each with 3 false successes; C2 29, C3 24, C0 22 with
none; 10 or 11 of the 12 tasks solved. A paired comparison of C4 with C0 gives
$p = 0.001953$.

\textbf{Budget:} Counted in runs and attempts only: 12 $\times$ 5 $\times$ 3
= 180 runs, four attempts per run in C1 and C2, one high-effort model. Tokens
and cost are not reported; they would be needed to compare the workflows at
equal spend.

\textbf{Variance:} Each cell was replicated three times, and iteration is said
to help consistency more than coverage, but no per-replicate spread or
confidence interval is in the range.

\textbf{Artifact:} The checker (CAPRI, with a full audit trail) and the small
author-maintained set of 12 tasks; nothing in the range says they are public.
\end{answer}

\end{document}
